\documentclass[12pt,a4paper]{article}

% ============================================================================
% PACKAGES
% ============================================================================
\usepackage[utf8]{inputenc}
\usepackage[T1]{fontenc}
\usepackage[french]{babel}
\usepackage{geometry}
\usepackage{graphicx}
\usepackage{xcolor}
\usepackage{listings}
\usepackage{hyperref}
\usepackage{booktabs}
\usepackage{array}
\usepackage{longtable}
\usepackage{fancyhdr}
\usepackage{titlesec}
\usepackage{tcolorbox}
\usepackage{fontawesome5}
\usepackage{enumitem}
\usepackage{caption}
\usepackage{subcaption}
\usepackage{float}
\usepackage{amsmath}
\usepackage{amssymb}

% ============================================================================
% CONFIGURATION
% ============================================================================
\geometry{margin=2.5cm}

% Couleurs personnalisées
\definecolor{v2vprimary}{RGB}{76, 175, 80}      % Vert V2V
\definecolor{v2vsecondary}{RGB}{43, 43, 43}     % Gris foncé
\definecolor{v2vlight}{RGB}{240, 248, 255}      % Bleu clair
\definecolor{codebackground}{RGB}{245, 245, 245}
\definecolor{codecomment}{RGB}{0, 128, 0}
\definecolor{codekeyword}{RGB}{0, 0, 255}
\definecolor{codestring}{RGB}{163, 21, 21}

% Configuration des liens
\hypersetup{
    colorlinks=true,
    linkcolor=v2vprimary,
    filecolor=v2vprimary,
    urlcolor=v2vprimary,
    citecolor=v2vprimary,
    pdftitle={Manuel d'Utilisation - V2V Simulator},
    pdfauthor={Équipe V2V},
    pdfsubject={Simulateur de Véhicules Communicants},
    pdfkeywords={V2V, simulation, VANET, communication véhiculaire}
}

% Configuration des listings (code)
\lstset{
    backgroundcolor=\color{codebackground},
    basicstyle=\ttfamily\small,
    breaklines=true,
    captionpos=b,
    commentstyle=\color{codecomment},
    frame=single,
    keywordstyle=\color{codekeyword}\bfseries,
    language=bash,
    numbers=left,
    numbersep=5pt,
    numberstyle=\tiny\color{gray},
    showstringspaces=false,
    stringstyle=\color{codestring},
    tabsize=4,
    xleftmargin=2em,
    framexleftmargin=1.5em
}

% Configuration des titres
\titleformat{\section}
    {\Large\bfseries\color{v2vprimary}}
    {\thesection}{1em}{}
\titleformat{\subsection}
    {\large\bfseries\color{v2vsecondary}}
    {\thesubsection}{1em}{}
\titleformat{\subsubsection}
    {\normalsize\bfseries\color{v2vsecondary}}
    {\thesubsubsection}{1em}{}

% En-têtes et pieds de page
\pagestyle{fancy}
\fancyhf{}
\fancyhead[L]{\textcolor{v2vprimary}{\textbf{V2V Simulator}}}
\fancyhead[R]{\textcolor{gray}{Manuel d'Utilisation}}
\fancyfoot[C]{\thepage}
\renewcommand{\headrulewidth}{1pt}
\renewcommand{\headrule}{\hbox to\headwidth{\color{v2vprimary}\leaders\hrule height \headrulewidth\hfill}}

% Boîtes personnalisées
\newtcolorbox{notebox}{
    colback=v2vlight,
    colframe=v2vprimary,
    title={\faInfoCircle\ Note},
    fonttitle=\bfseries,
    boxrule=1pt,
    arc=3pt
}

\newtcolorbox{warningbox}{
    colback=yellow!10,
    colframe=orange,
    title={\faExclamationTriangle\ Attention},
    fonttitle=\bfseries,
    boxrule=1pt,
    arc=3pt
}

\newtcolorbox{tipbox}{
    colback=green!5,
    colframe=v2vprimary,
    title={\faLightbulb\ Astuce},
    fonttitle=\bfseries,
    boxrule=1pt,
    arc=3pt
}

% ============================================================================
% DOCUMENT
% ============================================================================
\begin{document}

% ============================================================================
% PAGE DE TITRE
% ============================================================================
\begin{titlepage}
    \centering
    \vspace*{2cm}
    
    {\Huge\bfseries\textcolor{v2vprimary}{V2V Simulator}}\\[0.5cm]
    {\LARGE\textcolor{v2vsecondary}{Simulateur de Véhicules Communicants}}\\[2cm]
    
    \rule{\linewidth}{1pt}\\[0.5cm]
    {\huge\bfseries Manuel d'Utilisation}\\[0.5cm]
    \rule{\linewidth}{1pt}\\[2cm]
    
    {\Large Version 1.0}\\[0.5cm]
    {\large Décembre 2024}\\[3cm]
    
    \begin{center}
    \begin{tabular}{rl}
        \textbf{Technologies :} & C++20, Qt6, Boost.Graph, QXmlStreamReader \\[0.3cm]
        \textbf{Plateforme :} & Linux (Ubuntu 22.04+) \\[0.3cm]
        \textbf{Licence :} & Open Source \\
    \end{tabular}
    \end{center}
    
    \vfill
    
    {\large\textit{Projet de Réseaux Mobiles}}\\
    {\large Université de Haute-Alsace}
    
\end{titlepage}

% ============================================================================
% TABLE DES MATIÈRES
% ============================================================================
\tableofcontents
\newpage

% ============================================================================
% INTRODUCTION
% ============================================================================
\section{Introduction}

\subsection{Présentation du Projet}

Le \textbf{V2V Simulator} est un simulateur de communication véhicule-à-véhicule (Vehicle-to-Vehicle) développé en C++20 avec Qt6. Il permet de visualiser et d'analyser les communications entre véhicules autonomes dans un environnement urbain réaliste basé sur des données cartographiques OpenStreetMap.

\subsection{Objectifs}

\begin{itemize}[leftmargin=2cm]
    \item[\faCheck] Simuler le déplacement de véhicules sur un réseau routier réel
    \item[\faCheck] Modéliser les communications sans fil entre véhicules
    \item[\faCheck] Visualiser les graphes d'interférences et de connectivité
    \item[\faCheck] Analyser les performances du réseau V2V (VANET)
    \item[\faCheck] Étudier l'impact de la densité de véhicules sur la connectivité
\end{itemize}

\subsection{Fonctionnalités Principales}

\begin{enumerate}
    \item \textbf{Visualisation Interactive} : Carte OpenStreetMap avec zoom et panoramique
    \item \textbf{Simulation Multi-Véhicules} : Jusqu'à 5000 véhicules simultanés
    \item \textbf{Graphe Routier} : Navigation réaliste sur routes réelles
    \item \textbf{Communications V2V} : Modèle de propagation radio simplifié
    \item \textbf{Graphe d'Interférences} : Visualisation des connexions entre véhicules
\end{enumerate}

% ============================================================================
% PRÉREQUIS
% ============================================================================
\section{Prérequis et Installation}

\subsection{Configuration Matérielle Recommandée}

\begin{table}[H]
\centering
\begin{tabular}{lll}
\toprule
\textbf{Composant} & \textbf{Minimum} & \textbf{Recommandé} \\
\midrule
Processeur & Intel Core i5 / AMD Ryzen 5 & Intel Core i7 / AMD Ryzen 7 \\
Mémoire RAM & 8 Go & 16 Go \\
Stockage & 2 Go disponibles & 5 Go SSD \\
Carte graphique & Intégrée & GPU dédié avec OpenGL 3.3+ \\
Résolution écran & 1920×1080 & 2560×1440 ou supérieur \\
\bottomrule
\end{tabular}
\caption{Configuration matérielle}
\end{table}

\subsection{Dépendances Logicielles}

\subsubsection{Système d'Exploitation}
\begin{itemize}
    \item Linux (Ubuntu 22.04 LTS ou supérieur recommandé)
    \item Kernel 5.15 ou supérieur
\end{itemize}

\subsubsection{Compilateur et Outils de Build}
\begin{lstlisting}[language=bash]
# Installation des outils de base
sudo apt update
sudo apt install -y build-essential cmake ninja-build git
sudo apt install -y g++-12  # Support C++20
\end{lstlisting}

\subsubsection{Bibliothèques Requises}

\begin{table}[H]
\centering
\begin{tabular}{llp{7cm}}
\toprule
\textbf{Bibliothèque} & \textbf{Version} & \textbf{Description} \\
\midrule
Qt6 & 6.4+ & Framework d'interface graphique et parser XML (QXmlStreamReader) \\
Boost & 1.74+ & Boost.Graph pour les algorithmes de graphe \\
TBB & 2021.5+ & Parallélisation Intel Threading Building Blocks \\
Eigen3 & 3.4+ & Bibliothèque d'algèbre linéaire \\
libcurl & 7.81+ & Téléchargement des tuiles OSM \\
SQLite3 & 3.37+ & Cache des tuiles de carte \\
\bottomrule
\end{tabular}
\caption{Bibliothèques requises}
\end{table}

\begin{lstlisting}[language=bash, caption={Installation des dépendances}]
# Qt6
sudo apt install -y qt6-base-dev qt6-tools-dev libqt6svg6-dev \
    libqt6opengl6-dev qt6-positioning-dev

# Boost
sudo apt install -y libboost-all-dev

# Autres bibliothèques
sudo apt install -y libtbb-dev libeigen3-dev libcurl4-openssl-dev \
    libsqlite3-dev libgl-dev
\end{lstlisting}

\subsection{Installation depuis les Sources}

\subsubsection{Clonage du Dépôt}
\begin{lstlisting}[language=bash]
git clone https://github.com/maliko18/v2v-simulation.git
cd v2v-simulation
\end{lstlisting}

\subsubsection{Compilation}
\begin{lstlisting}[language=bash]
# Créer le répertoire de build
mkdir cmake-build-debug && cd cmake-build-debug

# Configuration CMake
cmake -DCMAKE_BUILD_TYPE=Release -GNinja ..

# Compilation
ninja

# Lancement
./v2v_simulator
\end{lstlisting}

\begin{tipbox}
Pour une compilation optimisée, utilisez \texttt{-DCMAKE\_BUILD\_TYPE=Release} et ajoutez les flags \texttt{-O3 -march=native} pour des performances maximales.
\end{tipbox}

% ============================================================================
% INTERFACE UTILISATEUR
% ============================================================================
\section{Interface Utilisateur}

\subsection{Vue d'Ensemble}

L'interface du V2V Simulator est organisée en plusieurs zones fonctionnelles :

\begin{figure}[H]
\centering
\begin{tcolorbox}[colback=white, colframe=v2vsecondary, width=\textwidth]
\begin{verbatim}
┌─────────────────────────────────────────────────────────────────┐
│                         BARRE DE MENU                          │
├──────────────┬──────────────────────────────────────────────────┤
│              │                                                  │
│   PANNEAU    │                                                  │
│   DE         │              ZONE DE VISUALISATION               │
│   CONTRÔLE   │                  (MapView)                       │
│              │                                                  │
│  - Start     │         Carte OpenStreetMap                      │
│  - Pause     │         Véhicules (rectangles)                   │
│  - Reset     │         Connexions V2V (lignes vertes)           │
│  - Vitesse   │         Rayons de transmission (cercles bleus)   │
│  - Véhicules │                                                  │
│              │                                                  │
├──────────────┴──────────────────────────────────────────────────┤
│ Vehicles: 500 | Connections: 1234 | Time: 45.3s                │
└─────────────────────────────────────────────────────────────────┘
\end{verbatim}
\end{tcolorbox}
\caption{Structure de l'interface utilisateur}
\end{figure}

\subsection{Panneau de Contrôle (Gauche)}

Le panneau de contrôle situé à gauche de l'écran permet de configurer et piloter la simulation.

\subsubsection{Contrôles de Simulation}

\begin{table}[H]
\centering
\begin{tabular}{lp{10cm}}
\toprule
\textbf{Bouton} & \textbf{Fonction} \\
\midrule
\textcolor{v2vprimary}{\textbf{▶ Start}} & Démarre la simulation. Crée les véhicules si c'est le premier lancement. \\
\textbf{⏸ Pause} & Met en pause la simulation. Les véhicules restent à leur position actuelle. \\
\textbf{↻ Reset} & Arrête la simulation et supprime tous les véhicules. Remet le temps à zéro. \\
\bottomrule
\end{tabular}
\caption{Boutons de contrôle de simulation}
\end{table}

\subsubsection{Paramètres de Simulation}

\begin{table}[H]
\centering
\begin{tabular}{llp{7cm}}
\toprule
\textbf{Paramètre} & \textbf{Plage} & \textbf{Description} \\
\midrule
Speed (Vitesse) & 0.1x - 10.0x & Échelle de temps de la simulation \\
Vehicles (Véhicules) & 10 - 5000 & Nombre de véhicules à simuler \\
\bottomrule
\end{tabular}
\caption{Paramètres configurables}
\end{table}

\begin{notebox}
Le nombre de véhicules ne peut être modifié que lorsque la simulation est arrêtée. Cliquez sur \textbf{Reset} avant de changer ce paramètre.
\end{notebox}

\subsection{Zone de Visualisation (MapView)}

La zone centrale affiche la carte interactive avec les éléments de simulation.

\subsubsection{Éléments Visuels}

\begin{table}[H]
\centering
\begin{tabular}{lp{10cm}}
\toprule
\textbf{Élément} & \textbf{Description} \\
\midrule
\textbf{Fond de carte} & Tuiles OpenStreetMap téléchargées automatiquement et mises en cache \\
\textbf{Véhicules} & Rectangles colorés (noir, gris, blanc, bleu nuit) représentant chaque véhicule \\
\textbf{Connexions V2V} & Lignes vertes entre véhicules pouvant communiquer \\
\textbf{Rayons de transmission} & Cercles bleus semi-transparents autour de chaque véhicule \\
\textbf{Véhicule sélectionné} & Bordure jaune épaisse autour du véhicule sélectionné \\
\bottomrule
\end{tabular}
\caption{Éléments visuels de la carte}
\end{table}

\subsubsection{Couleurs des Véhicules}

Les véhicules sont représentés par des rectangles de 4 couleurs différentes, attribuées aléatoirement :

\begin{itemize}
    \item \textbf{Noir} : \texttt{\#000000}
    \item \textbf{Gris} : \texttt{\#808080}
    \item \textbf{Blanc} : \texttt{\#FFFFFF}
    \item \textbf{Bleu Nuit} : \texttt{\#191970} (Midnight Blue)
\end{itemize}

\subsubsection{Rayons de Transmission}

Chaque véhicule possède un rayon de transmission aléatoire parmi les valeurs suivantes :

\begin{itemize}
    \item 100 mètres
    \item 200 mètres
    \item 300 mètres
    \item 400 mètres
    \item 500 mètres
\end{itemize}

\begin{tipbox}
Le rayon de transmission détermine la portée des communications V2V. Deux véhicules sont connectés si la distance entre eux est inférieure au rayon maximal des deux véhicules.
\end{tipbox}

\subsection{Barre de Statut (Bas)}

La barre de statut affiche en temps réel les métriques de la simulation :

\begin{itemize}
    \item \textbf{Vehicles:} Nombre de véhicules actifs dans la simulation
    \item \textbf{Connections:} Nombre de liens V2V établis
    \item \textbf{Time:} Temps de simulation écoulé (en secondes)
\end{itemize}

\subsection{Panneau d'Information du Véhicule}

Lorsqu'un véhicule est sélectionné, un panneau d'information apparaît en bas à droite de l'écran :

\begin{tcolorbox}[colback=white, colframe=gray, width=0.5\textwidth]
\begin{verbatim}
Selected Vehicle #42
─────────────────────
Position: (47.7456, 7.3389)
Speed: 54.2 km/h
Direction: 125.3°
TX Radius: 300 m
Connections: 5
\end{verbatim}
\end{tcolorbox}

% ============================================================================
% CONTRÔLES ET NAVIGATION
% ============================================================================
\section{Contrôles et Navigation}

\subsection{Contrôles Souris}

\begin{table}[H]
\centering
\begin{tabular}{lp{10cm}}
\toprule
\textbf{Action} & \textbf{Effet} \\
\midrule
\textbf{Clic gauche + Glisser} & Déplacement (pan) de la carte \\
\textbf{Molette vers le haut} & Zoom avant (agrandissement) \\
\textbf{Molette vers le bas} & Zoom arrière (réduction) \\
\textbf{Clic gauche sur véhicule} & Sélection du véhicule \\
\bottomrule
\end{tabular}
\caption{Contrôles souris}
\end{table}

\subsection{Raccourcis Clavier}

\begin{table}[H]
\centering
\begin{tabular}{llp{8cm}}
\toprule
\textbf{Touche} & \textbf{Action} & \textbf{Description} \\
\midrule
\textbf{+} ou \textbf{=} & Zoom avant & Augmente le niveau de zoom \\
\textbf{-} ou \textbf{\_} & Zoom arrière & Diminue le niveau de zoom \\
\textbf{↑} (Flèche haut) & Déplacer vers le haut & Déplace la vue vers le nord \\
\textbf{↓} (Flèche bas) & Déplacer vers le bas & Déplace la vue vers le sud \\
\textbf{←} (Flèche gauche) & Véhicule précédent & Sélectionne le véhicule le plus proche à gauche \\
\textbf{→} (Flèche droite) & Véhicule suivant & Sélectionne le véhicule le plus proche à droite \\
\textbf{H} & Retour maison & Recentre la carte sur Mulhouse (position initiale) \\
\textbf{V} & Toggle véhicules & Affiche/masque les véhicules \\
\textbf{C} & Toggle connexions & Affiche/masque les connexions V2V \\
\textbf{R} & Toggle routes & Affiche/masque le graphe routier \\
\textbf{T} & Toggle rayons TX & Affiche/masque les rayons de transmission \\
\textbf{Échap} & Désélectionner & Annule la sélection du véhicule courant \\
\bottomrule
\end{tabular}
\caption{Raccourcis clavier}
\end{table}

\subsection{Navigation entre Véhicules}

La navigation par clavier permet de parcourir les véhicules de manière intuitive :

\begin{enumerate}
    \item \textbf{Flèche gauche (←)} : Sélectionne le véhicule le plus proche situé à \textbf{gauche} du véhicule actuel
    \item \textbf{Flèche droite (→)} : Sélectionne le véhicule le plus proche situé à \textbf{droite} du véhicule actuel
\end{enumerate}

\begin{warningbox}
Si aucun véhicule n'est sélectionné, les touches gauche/droite sélectionnent le premier véhicule visible à l'écran.
\end{warningbox}

% ============================================================================
% FONCTIONNEMENT DE LA SIMULATION
% ============================================================================
\section{Fonctionnement de la Simulation}

\subsection{Cycle de Simulation}

La simulation fonctionne selon le cycle suivant :

\begin{enumerate}
    \item \textbf{Initialisation} : Création des véhicules à des positions aléatoires sur le réseau routier
    \item \textbf{Mise à jour des positions} : Chaque véhicule se déplace selon sa vitesse et direction
    \item \textbf{Calcul des connexions} : Le graphe d'interférences est mis à jour périodiquement
    \item \textbf{Rendu graphique} : L'affichage est rafraîchi à 30 FPS
\end{enumerate}

\subsection{Modèle de Véhicule}

Chaque véhicule possède les attributs suivants :

\begin{table}[H]
\centering
\begin{tabular}{llp{7cm}}
\toprule
\textbf{Attribut} & \textbf{Type} & \textbf{Description} \\
\midrule
ID & Entier & Identifiant unique du véhicule \\
Position & (lat, lon) & Coordonnées GPS en degrés décimaux \\
Vitesse & m/s & Vitesse de déplacement (10-25 m/s, soit 36-90 km/h) \\
Direction & Radians & Angle de déplacement (0-2$\pi$) \\
Rayon TX & Mètres & Rayon de transmission radio (100-500 m) \\
Actif & Booléen & État du véhicule dans la simulation \\
\bottomrule
\end{tabular}
\caption{Attributs d'un véhicule}
\end{table}

\subsection{Modèle de Communication V2V}

\subsubsection{Critère de Connexion}

Deux véhicules $V_1$ et $V_2$ sont considérés comme \textbf{connectés} si :

\begin{equation}
d(V_1, V_2) \leq \max(R_1, R_2)
\end{equation}

Où :
\begin{itemize}
    \item $d(V_1, V_2)$ est la distance euclidienne entre les deux véhicules (en mètres)
    \item $R_1$ est le rayon de transmission de $V_1$
    \item $R_2$ est le rayon de transmission de $V_2$
\end{itemize}

\subsubsection{Calcul de Distance}

La distance entre deux véhicules est calculée en utilisant une approximation plane :

\begin{equation}
d = \sqrt{(\Delta lat \times 111320)^2 + (\Delta lon \times 111320)^2}
\end{equation}

Où 111320 mètres correspond approximativement à 1 degré de latitude/longitude à nos latitudes.

\subsection{Graphe d'Interférences}

Le graphe d'interférences est une structure de données qui maintient les connexions entre véhicules :

\begin{itemize}
    \item \textbf{Nœuds} : Chaque véhicule est un nœud du graphe
    \item \textbf{Arêtes} : Une arête existe entre deux véhicules s'ils peuvent communiquer
    \item \textbf{Mise à jour} : Le graphe est recalculé périodiquement (toutes les 15-30 frames selon le nombre de véhicules)
\end{itemize}

\begin{notebox}
Pour optimiser les performances, la fréquence de mise à jour du graphe d'interférences s'adapte automatiquement au nombre de véhicules :
\begin{itemize}
    \item < 500 véhicules : mise à jour toutes les 15 frames
    \item 500-1000 véhicules : mise à jour toutes les 15 frames
    \item 1000-1500 véhicules : mise à jour toutes les 20 frames
    \item > 1500 véhicules : mise à jour toutes les 30 frames
\end{itemize}
\end{notebox}

% ============================================================================
% ARCHITECTURE TECHNIQUE
% ============================================================================
\section{Architecture Technique}

\subsection{Structure du Projet}

\begin{lstlisting}[language=bash, caption={Arborescence du projet}]
v2v-simulator/
├── include/              # Fichiers d'en-tête (.hpp)
│   ├── core/             # Vehicle, SimulationEngine
│   ├── network/          # RoadGraph, InterferenceGraph, PathPlanner
│   ├── visualization/    # MainWindow, MapView
│   ├── data/             # OSMParser, TileManager, GeometryUtils
│   └── utils/            # Logger
├── src/                  # Implémentations (.cpp)
├── data/                 # Données OSM (Mulhouse, Alsace)
├── docs/                 # Documentation
├── osm_cache/            # Cache des tuiles OSM (généré)
├── CMakeLists.txt        # Configuration CMake
└── run_optimized.sh      # Script de lancement
\end{lstlisting}

\subsection{Modules Principaux}

\begin{table}[H]
\centering
\begin{tabular}{lll}
\toprule
\textbf{Module} & \textbf{Composants} & \textbf{Technologies} \\
\midrule
Core & Vehicle, SimulationEngine & Qt6, C++20 \\
Network & RoadGraph, InterferenceGraph, PathPlanner & Boost.Graph, R-tree \\
Visualization & MainWindow, MapView & Qt6, QPainter \\
Data & OSMParser, TileManager, GeometryUtils & QXmlStreamReader, CURL, SQLite \\
Utils & Logger & Qt6 \\
\bottomrule
\end{tabular}
\caption{Modules du système}
\end{table}

\subsection{Diagramme de Classes Simplifié}

\begin{verbatim}
                    ┌─────────────────┐
                    │   MainWindow    │
                    │  (QMainWindow)  │
                    └────────┬────────┘
                             │
              ┌──────────────┼──────────────┐
              ▼              ▼              ▼
     ┌────────────┐  ┌──────────────┐  ┌──────────┐
     │  MapView   │  │ Simulation   │  │  Status  │
     │ (QWidget)  │  │   Engine     │  │   Bar    │
     └────────────┘  └──────┬───────┘  └──────────┘
                            │
         ┌──────────────────┼──────────────────┐
         ▼                  ▼                  ▼
   ┌───────────┐     ┌─────────────┐    ┌───────────┐
   │ RoadGraph │     │ Interference│    │  Vehicle  │
   │           │     │   Graph     │    │           │
   └───────────┘     └─────────────┘    └───────────┘
\end{verbatim}

% ============================================================================
% PERFORMANCES ET OPTIMISATION
% ============================================================================
\section{Performances et Optimisation}

\subsection{Optimisations Implémentées}

\begin{enumerate}
    \item \textbf{Frustum Culling} : Seuls les véhicules visibles à l'écran sont dessinés
    \item \textbf{Mise à jour adaptative} : La fréquence de calcul du graphe d'interférences s'adapte à la charge
    \item \textbf{Cache de tuiles} : Les tuiles OSM sont mises en cache localement (SQLite)
    \item \textbf{Limitation du rendu} : Maximum 2000 véhicules affichés simultanément
    \item \textbf{Limitation des connexions} : Maximum 1000 lignes de connexion dessinées
\end{enumerate}

\subsection{Recommandations d'Utilisation}

\begin{table}[H]
\centering
\begin{tabular}{lll}
\toprule
\textbf{Nombre de véhicules} & \textbf{Performance attendue} & \textbf{Recommandation} \\
\midrule
< 500 & Excellente (60 FPS) & Toutes options activées \\
500 - 1000 & Bonne (30-60 FPS) & Connexions activables \\
1000 - 2000 & Acceptable (20-30 FPS) & Désactiver rayons TX \\
> 2000 & Variable & Désactiver connexions \\
\bottomrule
\end{tabular}
\caption{Performances selon le nombre de véhicules}
\end{table}

\begin{warningbox}
Au-delà de 2000 véhicules, il est recommandé de désactiver l'affichage des connexions (touche \textbf{C}) et des rayons de transmission (touche \textbf{T}) pour maintenir des performances acceptables.
\end{warningbox}

% ============================================================================
% DÉPANNAGE
% ============================================================================
\section{Dépannage}

\subsection{Problèmes Courants}

\begin{longtable}{p{5cm}p{9cm}}
\toprule
\textbf{Problème} & \textbf{Solution} \\
\midrule
\endhead
La carte n'affiche pas de tuiles & Vérifiez votre connexion Internet. Les tuiles sont téléchargées depuis les serveurs OpenStreetMap. Vérifiez aussi que le dossier \texttt{osm\_cache} est accessible en écriture. \\
\midrule
Les véhicules n'apparaissent pas & Appuyez sur \textbf{V} pour activer l'affichage des véhicules. Vérifiez que la simulation est démarrée (bouton Start). \\
\midrule
L'application est lente & Réduisez le nombre de véhicules. Désactivez les connexions (C) et les rayons (T). \\
\midrule
Compilation échoue & Vérifiez que toutes les dépendances sont installées. Consultez les messages d'erreur CMake. \\
\midrule
Pas de connexions visibles & Appuyez sur \textbf{C} pour activer l'affichage des connexions. Les connexions ne s'affichent que si moins de 300 véhicules sont visibles. \\
\midrule
Le zoom ne fonctionne pas & Cliquez sur la carte pour lui donner le focus, puis utilisez la molette. \\
\bottomrule
\caption{Guide de dépannage}
\end{longtable}

\subsection{Fichiers de Log}

Les logs de l'application sont enregistrés dans le fichier \texttt{v2v\_simulator.log} situé dans le répertoire d'exécution.

\begin{lstlisting}[language=bash, caption={Consultation des logs}]
# Afficher les dernières lignes du log
tail -f v2v_simulator.log

# Rechercher des erreurs
grep -i error v2v_simulator.log
\end{lstlisting}

% ============================================================================
% ANNEXES
% ============================================================================
\section{Annexes}

\subsection{Formule de Haversine}

Pour les calculs de distance géographique précis, la formule de Haversine est utilisée :

\begin{equation}
d = 2R \cdot \arcsin\left(\sqrt{\sin^2\left(\frac{\phi_2-\phi_1}{2}\right) + \cos(\phi_1)\cos(\phi_2)\sin^2\left(\frac{\lambda_2-\lambda_1}{2}\right)}\right)
\end{equation}

Où :
\begin{itemize}
    \item $R$ = rayon de la Terre (6371 km)
    \item $\phi_1, \phi_2$ = latitudes en radians
    \item $\lambda_1, \lambda_2$ = longitudes en radians
\end{itemize}

\subsection{Conversion Coordonnées}

\subsubsection{Latitude/Longitude vers Pixels (Web Mercator)}

\begin{equation}
x = \frac{(\lambda + 180)}{360} \times 2^{zoom} \times 256
\end{equation}

\begin{equation}
y = \left(1 - \frac{\ln(\tan(\phi) + \sec(\phi))}{\pi}\right) \times \frac{2^{zoom} \times 256}{2}
\end{equation}

\subsection{Références}

\begin{enumerate}
    \item OpenStreetMap : \url{https://www.openstreetmap.org}
    \item Qt Documentation : \url{https://doc.qt.io/qt-6/}
    \item Boost.Graph : \url{https://www.boost.org/doc/libs/release/libs/graph/}
    \item IEEE 802.11p (WAVE) : Standard pour les communications V2V
    \item ETSI ITS-G5 : Standard européen pour les systèmes de transport intelligents
\end{enumerate}

% ============================================================================
% CONCLUSION
% ============================================================================
\section*{Conclusion}
\addcontentsline{toc}{section}{Conclusion}

Le V2V Simulator offre une plateforme complète pour l'étude et la visualisation des communications véhicule-à-véhicule. Son interface intuitive, combinée à des performances optimisées, permet d'explorer différents scénarios de densité de véhicules et d'analyser leur impact sur la connectivité du réseau.

Pour toute question ou contribution, n'hésitez pas à ouvrir une issue sur le dépôt GitHub :

\begin{center}
\url{https://github.com/maliko18/v2v-simulation}
\end{center}

\vspace{1cm}
\begin{center}
\textcolor{v2vprimary}{\Large\textbf{Bon usage du V2V Simulator !}}
\end{center}

\end{document}
